\documentclass[10pt,a4paper]{amsart}
\usepackage[margin=19mm]{geometry}
\usepackage{amsmath,amssymb,mathtools}
\usepackage[scr=rsfs,cal=euler]{mathalfa}
\usepackage{xfrac}
\usepackage[unicode,hidelinks,bookmarks=false]{hyperref}
\newcommand{\ie}{i.e.\ }
\newcommand{\cf}{cf.\ }
\newcommand{\resp}{respectively\ }
\newcommand{\Subsection}{\subsection}
\providecommand{\nobreakdash}{\nobreak\hbox{-}}
\setlength{\emergencystretch}{2em}
\multlinegap0pt
\usepackage{amssymb}
\usepackage{mathtools}
\usepackage[shortlabels]{enumitem}
\usepackage{tikz}
\newtheorem{lemma}{Lemma}
\usepackage{cleveref}
\crefname{lemma}{Lemma}{Lemmas}
\newcommand{\D}{\mathcal D}
\newcommand{\King}{\mathcal K}
\title{The Multiset Dimension of Graphs:\\ Extremal Values and King Grids}
\author{Jaan Allikvere}
\date{Source: arXiv:2607.28813v1 (CC BY 4.0)}
\begin{document}
\maketitle

\noindent\emph{Source and licence.} The Multiset Dimension of Graphs:\\ Extremal Values and King Grids. Version arXiv:2607.28813v1 (CC BY 4.0), distributed by arXiv under the Creative Commons Attribution 4.0 International licence, http://creativecommons.org/licenses/by/4.0/, as recorded in the arXiv licence metadata for this exact version and shown on the abstract page https://arxiv.org/abs/2607.28813v1. Full licence notice: \texttt{licenses/arxiv-2607.28813-CC-BY-4.0.txt}.\par\smallskip

\noindent\textit{Editorial scope.} Counted unit: the article's lemma lem:distinct (source0.tex lines 510--512). The article's complete original proof of it is delivered with it (source0.tex lines 514--620, one \texttt{proof} environment of 107 lines). No mathematics is written, repaired or re-derived: the statement and the proof are the article's own lines, in the article's own notation, and no retained line is substituted or re-flowed.  Retained uncounted as the source-local context the proof needs: lines 395--399; lines 429--439; lines 431--438; lines 470--478.  Nothing was omitted from any retained span, so the package carries no omission record.  No retained line contains a citation, so the wrapper carries no bibliography and no bibliography entry.  No retained line contains a figure environment, an image-inclusion command or drawing code, so no image is delivered and the package declares no asset dependency.  Editorial formatting only, in addition: an AMS \texttt{amsart} preamble that restates the article's own declarations and macros used by the retained lines, namely the environments \texttt{lemma} on separate counters, together with the standard packages those lines need.  The retained environments print in this document as Lemma 1 in order of appearance, whereas the article prints the same items with the numbering of its own full document.  The complete native source of the article as distributed in the arXiv e-print for this exact version is bundled unchanged as \texttt{sources/206442/source0.tex}.
\begin{equation}\label{eq:inverse}
 x=\frac{u+v+m}{2},
 \qquad
 y=\frac{u-v+m}{2}.
\end{equation}

\begin{lemma}[Boundary inequalities]\label{lem:boundary}
If $S$ is an ordinary resolving set of $\King_m$, then
\begin{equation}\label{eq:boundaryineq}
\begin{aligned}
 u_-+v_-&\le -m, &
 u_- -v_+&\le -m,\\
 u_++v_+&\ge m, &
 u_+-v_-&\ge m.
\end{aligned}
\end{equation}
\end{lemma}

\begin{equation}
\begin{aligned}
 u_-+v_-&\le -m, &
 u_- -v_+&\le -m,\\
 u_++v_+&\ge m, &
 u_+-v_-&\ge m.
\end{aligned}
\end{equation}

\begin{lemma}[Antipodal obstruction]\label{lem:antipodal}
Suppose two landmarks in $\D_m$ are antipodal, say
\[
 A=(p,q),
 \qquad
 C=(-p,-q).
\]
Then no choice of a third landmark makes $\{A,B,C\}$ a multiset resolving set.
\end{lemma}

\begin{lemma}\label{lem:distinct}
If all three $u$-coordinates and all three $v$-coordinates are distinct, then the three landmarks do not form a multiset resolving set.
\end{lemma}

\begin{proof}
Suppose, for a contradiction, that the three landmarks form a multiset resolving set.  In particular they form an ordinary resolving set, so the boundary inequalities of \cref{lem:boundary} hold.

Suppose first that
\[
 v_1<v_2<v_3.
\]
By \cref{lem:boundary},
\[
 u_1+v_1\le -m,
 \qquad
 u_3+v_3\ge m.
\]
On the other hand, membership in $\D_m$ gives the reverse inequalities.  Hence
\[
 u_1+v_1=-m,
 \qquad
 u_3+v_3=m.
\]
The remaining two boundary inequalities become
\[
 u_1-v_3\le -m,
 \qquad
 u_3-v_1\ge m.
\]
Substituting the two equalities above gives $v_1+v_3\ge0$ and $v_1+v_3\le0$.  Therefore
\[
 (u_3,v_3)=-(u_1,v_1),
\]
and \cref{lem:antipodal} applies.  The decreasing monotone order is equivalent by reflection.

It remains to consider the zigzag orbit.  By symmetry, assume
\begin{equation}\label{eq:zigzagorder}
 v_3<v_1<v_2.
\end{equation}
The last inequality in \eqref{eq:boundaryineq}, together with $p_3\in\D_m$, forces
\begin{equation}\label{eq:p3edge}
 u_3-v_3=m.
\end{equation}
Indeed, $u_3-v_3\ge m$, while $u_3-v_3\le |u_3|+|v_3|\le m$.  In particular $u_3\ge0\ge v_3$ and $v_3=u_3-m$.  Moreover $u_3\ge1$: if $u_3=0$, then $u_+=0$, and the third boundary inequality gives $v_+=v_2=m$, which forces $u_2=0=u_3$, contradicting the strict ordering of the $u_i$.

If $u_2\le u_3-2$, then the two vertices
\[
 P=(u_3-2,v_3),
 \qquad
 Q=(u_3-1,v_3+1)
\]
belong to $\D_m$ and have equal ordered distance vectors to all three landmarks.  Membership uses $u_3\ge1$: since $v_3=u_3-m$, we get $|u_3-2|+|v_3|=m-2$ for $u_3\ge2$ and $=m$ for $u_3=1$, while $|u_3-1|+|v_3+1|=m-2$ for $u_3\le m-1$ and $=m$ for $u_3=m$; both points inherit the parity of $p_3$.  For $i=1,2$, the $u$-distance increases by one from $P$ to $Q$ while the $v$-distance decreases by one; for $i=3$, both $L^1$-distances equal $2$.  Hence ordinary resolvability requires
\begin{equation}\label{eq:uadj}
 u_2=u_3-1.
\end{equation}
We next show that ordinary resolvability also requires
\begin{equation}\label{eq:vadj}
 v_1=v_3+1.
\end{equation}
Indeed, suppose that $v_1\ge v_3+2$.  In view of \eqref{eq:uadj}, the two vertices
\[
 P'=(u_3,v_3+2),
 \qquad
 Q'=(u_3-1,v_3+1)
\]
belong to $\D_m$.  To see this, note first that $v_3\ne0$: otherwise $u_3=m$, and the strict inequalities $0=v_3<v_1<v_2$ are incompatible with $p_2=(m-1,v_2)\in\D_m$.  Hence $v_3\le-1$.  If $v_3=-1$, then $P'=(m-1,1)$ lies on the boundary of the diamond; if $v_3\le-2$, then $|P'_u|+|P'_v|=m-2$.  In either case $P'\in\D_m$, while $|Q'_u|+|Q'_v|=m-2$; parity is preserved for both points.  For $i=1,2$, the $u$-distance decreases by one from $P'$ to $Q'$ while the $v$-distance increases by one, and both points have $L^1$-distance $2$ from $p_3$.  Thus $P'$ and $Q'$ have equal ordered distance vectors, a contradiction.  Since $v_1>v_3$, \eqref{eq:vadj} follows.

Write
\[
 v_3=-a-1,
 \qquad
 L=m-a.
\]
Then \eqref{eq:p3edge} gives $u_3=L-1$.  Relations \eqref{eq:uadj} and \eqref{eq:vadj}, the boundary inequalities, the diamond constraint, and the parity condition yield
\begin{equation}\label{eq:zigzagnormal}
\begin{aligned}
 p_1&=(-L,-a),\\
 p_2&=(L-2,a+2),\\
 p_3&=(L-1,-a-1),
\end{aligned}
\qquad
0\le a\le m-2.
\end{equation}
For completeness, the only numerical choices left by the inequalities are
\[
 -L\le u_1\le -L+1,
 \qquad
 a+1\le v_2\le a+2;
\]
the parity condition selects $u_1=-L$ and $v_2=a+2$.

Returning to the original coordinates via \eqref{eq:inverse}, the landmarks are
\[
 (0,a),
 \qquad
 (m,m-a-2),
 \qquad
 (m-a-1,m).
\]
The distinct vertices
\[
 U=(m-a-2,m),
 \qquad
 V=(m-a,m)
\]
have the same ordered distance vector
\[
 (L,a+2,1).
\]
Thus the zigzag case is not even ordinarily resolving.
\end{proof}

\end{document}
